\chapter{Versionierte Eingaben und Ergebnisartefakte}
\label{api:artifacts}\index{Schema}\index{Provenienz}\index{VTU}
\begingroup\raggedright
\section{Verträge und Beispieldateien}
Ein Schema ist ein eigenständiger Vertrag. Paketversion 1.x, Model-Schema 1/2/3,
Bundle-Schema 1/2 und Zeitserien-Schema 1/2 sind nicht austauschbar.
Das Formatkapitel des Benutzerbands beschreibt vollständige TOML-/CSV-Rezepte;
hier werden die maschinenlesbaren Beziehungen und Reader-Grenzen festgelegt.
Änderungen der Signatur-/Typ-/Quellverträge erkennt das API-Inventar, die
Rezeptmatrix bindet Beispiele bytegenau an Produktionsinputs. Die Tests lesen
mit Model/Study/Surrogate-/Serienparsern, keinem eigens permissiven Testreader.

Model-Schema 1: SI-Deklaration, stationäre Physik, mesh.path/length\_scale\_to\_m,
Materialkatalog, Regionen, Randbedingungen/Lasten, Solveroptionen.
Schema 2 ergänzt Zeitprofile, CSV-Referenzen, Anfangs-/Ausgabe-/Probedaten für
Wärme. Schema 3 ergänzt mode/transient/harmonic und deren begrenzte
piezoelektrische Daten. Kein fremder Schlüssel, undefinierter Pfad, unpassende
Einheit oder ungültige Shape wird still angenommen. Ausführbare Quellen:
\nolinkurl{examples/thermopiezoelectric/case.toml},
\nolinkurl{examples/transient_heat/source.csv},
\nolinkurl{examples/piezo_dynamics/drive.csv}.

Study-Schema 1: relativ lokales Model und explizite parameter(name,target,unit,values).
Persistierter Katalog enthält aktuelle/retired Fälle, Versuchszustand,
Resultatpfade/-hashes, Solver-/Input-/Versionsfingerprints und Verlauf.
Quantity-Tabelle Schema 1 enthält columns, rows, catalog\_sha256 und
csv\_sha256; CSV-Spaltenreihenfolge wird aus dem Header gelesen, nicht aus
sortierter JSON-Schlüsselreihenfolge. Status, Fehlgründe, Balance-/Solver-
Verfügbarkeit und ungeprüfte Netzkonvergenz bleiben sichtbar.
Surrogat-Schema 1 enthält Parameter/Ziele, Basis, Trainingsskalierung,
Koeffizienten und Train/Test-Bericht samt Daten-/Versionsprovenienz.
Uncertainty-/Sampling-Schema 1 bindet explizite Gesetze, gemeinsame Zeilen,
RNG/Seed/Korrelation und numerische Budgets getrennt; siehe UQ-Abschnitt.

\section{Gmsh und Ergebnisbündel}
\texttt{read\_gmsh}: Gmsh 2.2/4.1 über meshio, explizite Länge-in-m-Skala,
Mehrfachgruppen in benannten Sätzen. Nur unterstützte 3D-Zellen/Flächen;
unsupported/geometrisch fehlerhafte Daten werden diagnostiziert.
\texttt{read\_vtu} ist ein erreichbarer Forschungsadapter, keine zusätzliche
stabile Datenmigrationszusage. \texttt{write\_vtu} nimmt point\_data[n,...]
und je Zelltyp cell\_data[e,...] sowie Satzmasken; korrekte Zelltypformen
und Felddefinitionen bleiben Aufgabe des Aufrufers.

\texttt{export\_result\_bundle(path, mesh, result)} liefert VTU-/JSON-Pfade,
überschreibt bestehende Ziele nicht und prüft Resultat-/Netzbindung.
Schema 1 beschreibt Tet4; Schema 2 erweitert um Tet10/Zelltyp-/Feldpositions-
metadaten. Primärfelder werden in Netzindexfolge mit Gültigkeitsmasken
exportiert; inaktive Punkte sind Padding, keine gelösten Werte.
Rohzellergebnisse bleiben in aktiver Volumenzellfolge, Facettenpads und
Schwerpunktwerte sind ausdrücklich markiert. JSON hält Metriken,
Feldbeschreibungen, SI-Einheiten, Komponenten und die VTU-Prüfsumme.
RunProvenance Schema 1 bindet Eingabehash, Solverparameter und Versionen;
es belegt effektive diskrete Inputs, keinen externen Callbackquelltext.
Hashidentität ersetzt keine physikalische Bilanz, Konvergenz oder Validierung.

\section{Zeitserien und komplexe Felder}
HeatSeriesWriter/\texttt{export\_heat\_series} schreiben netzgebundene
PVD-/VTU-Frames, Messpunkt-/Bilanz-CSV und Manifest. Piezo-Series benötigt
Schema-3-ModelResult/PiezoDynamicResult und erhält Moden, Newmark oder Harmonik.
Manifest Schema 1 Tet4 bzw. 2 Tet10 benennt collection coordinate ausdrücklich
als time\_s, frequency\_hz oder mode\_index. Komplexe u/v/a/phi werden vor
Amplitude-/Phasenauswertung interpoliert; Phasen rad, Spitzenamplituden SI.
Nullamplitude hat phase=0 und phase\_valid=0. Reale Moden/Zeitläufe erzeugen
keine fingierte Phase. Feld-/Zell-/Randmasken und Tet10-Schwerpunktmetadaten
bleiben prüfbar; Messpunkte sind Netzindizes oder physische Punkte[3] m.

Serienexport in neues Verzeichnis; abgebrochene Serien bleiben ausdrücklich
unvollständig mit ihren tatsächlich fertig geschriebenen Frames.
Ein vorhandenes Manifest ohne complete-Erfolg darf nicht als voller Lauf
interpretiert werden. Exporter ersetzen keine Messdatenerhebung und keine
fehlende Quellenlizenz. Reader/Tests:
\nolinkurl{tests/unit/test_result_bundle.py},
\nolinkurl{tests/unit/test_heat_series.py},
\nolinkurl{tests/unit/test_piezo_series.py},
\nolinkurl{tests/unit/test_study_reports.py},
\nolinkurl{tests/unit/test_surrogates.py}.

\section{Kompatibilität und Pflege}
Stabile Minor-Releases sind additiv. Inkompatible Einheiten, Vorzeichen,
Formen, Importpfade oder öffentliche Schemata verlangen die dokumentierte
Versions- und Migrationsentscheidung. Kein stilles Lesen eines neuen inkompatiblen
Schemas als altes. Deprecation braucht Warnung, dokumentierten Ersatz,
Migrationsbeispiel und mindestens ein Minor-Release Übergang; Entfernung
in Major. Unter \nolinkurl{Doc/09_RELEASE_PROCESS.md} stehen konkrete Gates.
Generierte Signaturen werden mit \nolinkurl{scripts/api_contracts.py}
aktualisiert, manuelle Fachverträge separat geprüft. Neue Exporte brauchen
Inventar-/Matrixeintrag und Beispiel; Produktions-API wird nicht nur zur
Vereinfachung der Dokumentation umgebaut.
\endgroup
